\section{Un flux aléatoire}
\label{sec:U4}

Jusqu'ici, les flux étaient certains : seule leur date variait. On ajoute
maintenant l'ingrédient qui fait basculer du monde bancaire au monde de
l'assurance --- le flux peut ne pas avoir lieu.

C'est l'étage décisif de la progression. Une fois franchi, la phrase qui
gouverne tout le dispositif prudentiel devient lisible : \emph{un contrat
d'assurance est une variable aléatoire}, et le provisionner consiste à en
calculer l'espérance actualisée.

\subsection{La notion : espérance actualisée d'un flux conditionnel}

\begin{notion}{Deux ingrédients qui se multiplient}
Soit un flux $C$ payable en $n$ si un événement se réalise, et rien sinon.
Notons $p$ la probabilité de cet événement. La valeur actuelle probable est
\[
  V_0 \;=\; \frac{\mathbb{E}[C_n]}{(1+r)^n} \;=\; \frac{p \times C}{(1+r)^n} .
\]
La formule de l'unité~1 est intacte ; il s'y ajoute un unique facteur.

\medskip
Ce facteur est d'une autre nature que le taux, et cette distinction structure
tout le reste du cours. La probabilité $p$ relève d'une \emph{hypothèse
actuarielle} --- mortalité, fréquence de sinistre, comportement de l'assuré ---
estimée sur des données d'expérience. Le taux $r$ relève d'une \emph{donnée de
marché}. Les deux se multiplient, mais ne se justifient pas, ne se révisent pas
et ne se contrôlent pas de la même manière.
\end{notion}

\begin{center}
\begin{tikzpicture}[x=1.5cm, y=1cm]
  \draw[timeaxis] (-0.3,0) -- (4.6,0);
  \foreach \x/\l in {0/0, 2/{\dots}, 4/n}
    { \node[timenode] at (\x,0) {}; \node[etiq, below=3pt] at (\x,0) {\l}; }
  \draw[flux] (4,0) -- (4,1.2);
  \node[etiq, above] at (4,1.25) {$C$};
  \node[etiq, right] at (4.05,0.65) {avec proba. $p$};
  \draw[dashed, grisactu] (4,0) -- (4,-0.8);
  \node[etiq, right] at (4.05,-0.55) {$0$ avec proba. $1-p$};
  \draw[fluxneg] (0,0.95) -- (0,0.05);
  \node[etiq, above] at (0,1.0) {$V_0 = p\,C\,v^{\,n}$};
\end{tikzpicture}
\end{center}

\begin{exercice}[capital sous condition de survie]
Un contrat prévoit le versement d'un capital de $50\,000$~\euro{} dans $8$~ans,
à la condition que l'assuré, âgé de $52$~ans, soit alors en vie. La probabilité
de survie sur cet horizon est de $0{,}942$. Le taux d'actualisation est de
$2{,}40\,\%$.

Calculer la valeur actuelle probable de cet engagement, puis la comparer à celle
du même capital versé de façon certaine.
\end{exercice}

\begin{corrige}
Avec $(1{,}024)^{8} = 1{,}20893$ :
\[
  V_0 = \frac{0{,}942 \times 50\,000}{1{,}20893}
      = \frac{47\,100}{1{,}20893} \approx 38\,960\ \text{\euro}.
\]
Le même capital versé de façon certaine vaudrait
$50\,000/1{,}20893 \approx 41\,359$~\euro.

\medskip
L'écart, $2\,399$~\euro, représente exactement $5{,}8\,\%$ de la valeur certaine
--- c'est-à-dire la probabilité de décès sur la période. L'aléa se paie donc en
proportion stricte de sa probabilité, ce qui n'ira plus de soi à l'unité~5.
\end{corrige}

\begin{exercice}[lecture inverse]
Un engagement de même nature --- capital de $50\,000$~\euro{} à $8$~ans, taux de
$2{,}40\,\%$ --- est provisionné pour $38\,960$~\euro. Quelle probabilité de
survie cette provision suppose-t-elle ?
\end{exercice}

\begin{corrige}
On inverse la relation :
\[
  p = \frac{V_0 \times (1+r)^{n}}{C}
    = \frac{38\,960 \times 1{,}20893}{50\,000}
    = \frac{47\,100}{50\,000} = 0{,}942 .
\]
La provision suppose une probabilité de survie de $94{,}2\,\%$.

Cette lecture inverse est un réflexe de contrôle utile : à partir d'un montant
provisionné, on peut toujours reconstituer l'hypothèse implicite et juger si
elle est plausible.
\end{corrige}

\subsection{Généralisation : rentes viagères et flux croissants}

\begin{generalisation}{Une pondération par période, quelle qu'en soit l'origine}
\textbf{La table de mortalité.} Une table donne les effectifs $l_x$ d'un groupe
initial encore en vie à l'âge $x$. On en déduit la probabilité qu'une personne
d'âge $x$ soit encore en vie $k$ années plus tard :
\[
  {}_{k}p_x \;=\; \frac{l_{x+k}}{l_x},
  \qquad\text{et la probabilité de décéder l'année } k :\quad
  \frac{l_{x+k-1}-l_{x+k}}{l_x}.
\]

\medskip
\textbf{La rente viagère.} Une rente de $1$ versée en fin d'année tant que
l'assuré est en vie, sur $n$ années au plus, vaut
\[
  a_{x:\overline{n}|} \;=\; \sum_{k=1}^{n} {}_{k}p_x \; v^{\,k}.
\]
C'est la rente certaine de l'unité~2 dont chaque terme est pondéré par sa
probabilité de survenance.

\medskip
\textbf{Les flux croissants.} Si le flux, au lieu d'être aléatoire, croît au
taux constant $g$ à l'infini, la même somme admet une forme fermée --- le modèle
de Gordon-Shapiro\index{Gordon-Shapiro} :
\[
  V_0 \;=\; \sum_{k=1}^{\infty} \frac{D_1 (1+g)^{k-1}}{(1+r)^{k}}
       \;=\; \frac{D_1}{r-g},
  \qquad\text{sous réserve que } g < r .
\]

\medskip
\textbf{Ce qu'il faut retenir.} Dans les deux cas, on applique au facteur
d'actualisation $v^{\,k}$ un \emph{multiplicateur propre à chaque période} :
une probabilité de survie ici, un facteur de croissance là. La mécanique est
identique ; seule l'interprétation du multiplicateur change.
\end{generalisation}

\begin{exercice}[rente viagère temporaire]
Un assuré de $65$~ans perçoit une rente de $6\,000$~\euro{} en fin de chacune
des cinq prochaines années, tant qu'il est en vie. Le taux est de $2{,}00\,\%$.
La table de mortalité donne :
\begin{center}\small
\begin{tabular}{lrrrrrr}
\toprule
$x$   & $65$ & $66$ & $67$ & $68$ & $69$ & $70$ \\
$l_x$ & $100\,000$ & $98\,600$ & $97\,100$ & $95\,400$ & $93\,500$ & $91\,400$ \\
\bottomrule
\end{tabular}
\end{center}
Calculer la valeur actuelle probable de la rente, puis la comparer à celle d'une
rente certaine de même montant et de même durée.
\end{exercice}

\begin{corrige}
\begin{center}\small
\begin{tabular}{crrrr}
\toprule
$k$ & ${}_{k}p_{65}$ & $v^{\,k}$ & ${}_{k}p_{65}\,v^{\,k}$ & Flux actualisé \\
\midrule
1 & $0{,}98600$ & $0{,}98039$ & $0{,}96667$ & $5\,800{,}00$ \\
2 & $0{,}97100$ & $0{,}96117$ & $0{,}93329$ & $5\,599{,}77$ \\
3 & $0{,}95400$ & $0{,}94232$ & $0{,}89898$ & $5\,393{,}85$ \\
4 & $0{,}93500$ & $0{,}92385$ & $0{,}86380$ & $5\,182{,}77$ \\
5 & $0{,}91400$ & $0{,}90573$ & $0{,}82784$ & $4\,967{,}03$ \\
\midrule
  & & & $\mathbf{4{,}49057}$ & $\mathbf{26\,943{,}42}$ \\
\bottomrule
\end{tabular}
\end{center}
La rente viagère vaut environ $26\,943$~\euro, contre
$6\,000 \times a_{\overline{5}|2\%} = 6\,000 \times 4{,}71346 \approx 28\,281$~\euro{}
pour une rente certaine.

\medskip
L'écart de $1\,337$~\euro{} est le seul effet de la mortalité. Noter que le
coefficient $a_{65:\overline{5}|} = 4{,}4906$ se lit directement : la rente
« vaut » $4{,}49$~années de versement en valeur d'aujourd'hui, contre $4{,}71$
si elle était certaine.
\end{corrige}

\begin{exercice}[capital décès et flux croissant]
\textbf{(a)} Sur la même table et au même taux, un contrat verse un capital de
$100\,000$~\euro{} en fin de l'année du décès, si celui-ci survient dans les
cinq prochaines années. Calculer la valeur actuelle probable de cet engagement.

\medskip
\textbf{(b)} Une action versera un dividende de $4{,}20$~\euro{} l'an prochain,
puis croissant de $2{,}50\,\%$ par an indéfiniment. Le taux de rendement exigé
est de $7{,}00\,\%$. Évaluer l'action, puis mesurer l'effet d'une révision de la
croissance à $3{,}00\,\%$.
\end{exercice}

\begin{corrige}
\textbf{(a)} La probabilité de décéder l'année $k$ se lit sur les différences
d'effectifs :
\begin{center}\small
\begin{tabular}{crrr}
\toprule
$k$ & Proba. de décès & $v^{\,k}$ & Flux actualisé \\
\midrule
1 & $0{,}01400$ & $0{,}98039$ & $1\,372{,}55$ \\
2 & $0{,}01500$ & $0{,}96117$ & $1\,441{,}75$ \\
3 & $0{,}01700$ & $0{,}94232$ & $1\,601{,}95$ \\
4 & $0{,}01900$ & $0{,}92385$ & $1\,755{,}31$ \\
5 & $0{,}02100$ & $0{,}90573$ & $1\,902{,}03$ \\
\midrule
  & $0{,}08600$ & & $\mathbf{8\,073{,}59}$ \\
\bottomrule
\end{tabular}
\end{center}
L'engagement vaut environ $8\,074$~\euro. Contrôle de cohérence : la somme des
probabilités de décès, $0{,}086$, est bien le complément de la probabilité de
survie à cinq ans, $0{,}914$.

\medskip
\textbf{(b)} Au taux de croissance initial :
\[
  V_0 = \frac{4{,}20}{0{,}07-0{,}025} = \frac{4{,}20}{0{,}045} = 93{,}33\ \text{\euro}.
\]
Avec $g = 3{,}00\,\%$ : $V_0 = 4{,}20/0{,}04 = 105{,}00$~\euro.

\medskip
Une révision de $50$~points de base sur la croissance déplace la valeur de
$12{,}5\,\%$. Cette instabilité n'est pas un défaut du modèle : elle est la
propriété de tout engagement de très longue durée, et elle explique pourquoi la
sensibilité d'une meilleure estimation à ses hypothèses fait l'objet d'une
attention réglementaire particulière.
\end{corrige}

\begin{pont}{Où cela sert dans Solvabilité~II}
Cette unité est la définition même de la meilleure estimation. Le texte de la
directive énonce que celle-ci est la moyenne des flux de trésorerie futurs
\emph{pondérés par leur probabilité}, compte tenu de la valeur temporelle de
l'argent. Les deux ingrédients de la notion, et rien d'autre.

\medskip
Trois conséquences opérationnelles en découlent.

\emph{Les hypothèses doivent être réalistes.} Un best estimate se calcule avec
des probabilités d'expérience, sans marge de prudence. C'est une différence de
principe avec le provisionnement en normes françaises, qui recourt à des tables
réglementaires délibérément prudentes. Les deux bilans peuvent donc afficher des
provisions très différentes pour un même contrat, sans qu'aucun ne soit erroné.

\emph{L'écart de provision a désormais deux sources.} L'unité~2 en avait
identifié une, l'actualisation. Celle-ci en apporte la seconde, le choix des
hypothèses. Ensemble, elles expliquent l'essentiel de la divergence entre bilan
comptable et bilan prudentiel.

\emph{La sensibilité aux hypothèses devient un objet en soi.} L'exercice~4.4
montre qu'un déplacement minime d'une hypothèse de long terme produit un effet
majeur sur la valeur.

\medskip
Il manque encore un ingrédient, et c'est celui qui rendra l'aléa non
proportionnel à sa probabilité : le contrat peut donner à l'assuré un
\emph{choix}. C'est l'objet de l'unité~5.
\end{pont}

\begin{refreg}
\textbf{Références.} Directive 2009/138/CE, article~77, paragraphe~2 : la
meilleure estimation correspond à la moyenne pondérée par leur probabilité des
flux de trésorerie futurs, compte tenu de la valeur temporelle de l'argent. \\
Le même article prévoit que le calcul repose sur des informations actualisées et
crédibles et sur des hypothèses réalistes. \\
Règlement délégué (UE) 2015/35 : hypothèses relatives aux options et garanties
contractuelles, et segmentation. \emph{[Références d'articles à préciser avant
diffusion.]}
\end{refreg}

\subsection*{Notations introduites}
\begin{description}[leftmargin=2.4cm, style=nextline, font=\normalfont]
  \notation{p}{probabilité de réalisation d'un flux}
  \notation{l_x}{effectif survivant à l'âge $x$ dans la table}
  \notation{{}_{k}p_x}{probabilité qu'une tête d'âge $x$ survive $k$ ans}
  \notation{a_{x:\overline{n}|}}{valeur actuelle d'une rente viagère temporaire de 1 versée pour $n$ années à une personne âgée de $x$ ans.}
  \notation{D_1}{premier dividende attendu}
  \notation{g}{taux de croissance des flux}
\end{description}
